\subsection{The last five-outside three-row act cones through nine states}

\begin{lemma}[The nine-state Nakayama calculation]
\label{C:Nakayama-eleven-nine-state}
Let $M$ have a $3\times2$ rectangular-band minimal ideal $K$ and
exactly five elements outside $K$.  Suppose its contracted rational
algebra is the Nakayama algebra
\[
 A=\mathbb QM/\mathbb QK,\qquad
 a\in e_1Ae_2,\quad b\in e_2Ae_1,\quad ab=c,\quad ba=0,
\]
and its two-dimensional row augmentation module is
$V=L=Ae_1/\mathbb Qc$.
There is no nine-element right $M$-set $S$ for which
$B\mathcal C(M,S)\simeq\mathbb CP^2$, with $SK$ consisting of
two two-element orbits.
\end{lemma}
\begin{proof}
We describe the complete finite calculation, its necessary algebraic
filters, and the actual cup certificates.  This is a calculation on
literal monoids and acts, not just on their linear representations.

First discard a split row augmentation module by
Proposition~\ref{C:split-row-act-cups}.  If $H^2(BM;\mathbb Q)=0$,
every degree-two class of the act cone is relative to $BM$; the
quotient by $BM$ is a suspension, so all its squares vanish.
Thus a possible source has a nonzero balanced form between its
row and column augmentation modules.  Since there are two columns,
such a form is a nonzero scalar-valued row functional.  Give the
three rows its affine values.  If all three values are distinct,
the finite affine self-maps of these values are the identity,
a possible common reflection, and constants.  Their augmentation
operators commute and diagonalize over $\mathbb Q$, so the row
module is split.  Otherwise relabeling and rescaling makes the
three values $(0,0,1)$.

Write an actor's row and column transformations as $\sigma$ and
$\tau$.  Compatibility with this functional is exactly
\[
 [\sigma(i)=2]-[\sigma(0)=2]
   =(\tau(1)-\tau(0))[i=2]\qquad(i=1,2).
\]
There are precisely $24$ pairs of transformations satisfying these
equations.  The generator
\path{code/a_eleven_three_by_two_targets.cpp} chooses every
multiset of four such pairs for the four nonidentity outside
elements: $\binom{27}{4}=17550$ profiles.  It discards $3550$
profiles whose row augmentation matrices commute and satisfy
$X^3=X$, which proves simultaneous splitting into rational
character lines.  For every remaining profile it tries every
outside product with the required row and column transformations,
and imposes all associativity equations.  Finally it takes orbits
under row, column, and outside-label permutations.  This gives
$120208$ associative sorted tables and $16399$ ordinary
isomorphism classes.  No passage to opposite monoids is made.
The source tables are retained in
\path{code/outputs/a_eleven_three_by_two_targets.tsv}.

The fourteen possible distinguished Nakayama bases have already
been independently classified in
Lemma~\ref{B:five-Nakayama-common-annihilator}.
The exact extraction
\path{code/a_nakayama11_sources.py} compares each contracted
table to all labelings of these bases.  It finds $93$ targets:
in $50$ the row module is the projective $Ae_2$, and in the
remaining $43$ it is $L$.  The projective case has zero
degree-two Tor and cannot give the fourth homology of a projective
plane.  Each of the $43$ records stores the full table, the
two rational characters, the row action, and exact distinguished-
basis coordinates of $e_1,b,c$.

Now suppose that a nine-state candidate is a projective plane.
Lemma~\ref{C:two-column-projective-orbit-kernel} identifies its
constant-action quotient with a compressed monoid whose column
augmentation module $W$ has dimension six and Tor dimensions
\[
 (\dim\operatorname{Tor}_0^A(W,L),
   \dim\operatorname{Tor}_1^A(W,L),
   \dim\operatorname{Tor}_2^A(W,L))=(2,0,1).
\]
The exact indecomposable table in
Lemma~\ref{B:five-exceptional-targets} forces
\[
 W=T_2\oplus e_1A\oplus e_2A.
\]
In particular $\dim We_1=\dim We_2=3$,
$\operatorname{rank}R_b=2$, and
$\operatorname{rank}R_c=1$.  If $\chi_1,\chi_2$ are the two
characters, the action on the seven-element quotient set must
therefore have
\[
 \#\operatorname{Fix}(m)=1+3\chi_1(m)+3\chi_2(m)
\]
fixed points for every actor $m$.  These are necessary conditions;
they are not used as substitutes for the final cup computation.

Here is the complete state normal form.  Label the core orbits
$\{0,1\}$ and $\{2,3\}$, with the second coordinate giving
the two columns, and label the other five states $4,\ldots,8$.
For each state $s$ let
$\gamma_s:\{0,1,2\}\to\{0,1\}$ record the core orbit reached
by the three rows.  A band element $(i,j)$ acts by
\[
 s(i,j)=2\gamma_s(i)+j.
\]
The two core orbits have constant $\gamma$ values zero and one.
Each outside actor preserves both core orbits and acts there by
its column transformation.  On the other states it must satisfy
\[
 \gamma_{sm}(i)=\gamma_s(\sigma_m(i)).
\]
Sort the five three-bit $\gamma$ values, retaining every map
within each equal-value fiber.  This gives $\binom{12}{5}=792$
profiles per source.  If all five values are constant, the act
splits into two invariant components and cannot have the required
first homology; discarding this case is valid.

The finite constraint search
\path{code/a_nakayama11_nine_states.cpp} assigns all twenty
unknown values of the four outside actors on the five remaining
states.  Its domains are exactly the displayed $\gamma$ fibers.
It propagates the fixed-point counts and every equation
$(sm)n=s(mn)$, branches on every value left in a domain, and
checks all $11^2\cdot9$ action equations on every completed act.
The $43\cdot792=34056$ profiles give $2119812$ completed acts
after the trace filter.  The ranks of $R_b,R_c$ leave $131144$
records.  Their six-by-six integer matrices have entries of
absolute value at most five.  Their ranks computed modulo
$1000000007$ are exact rational ranks: every minor is bounded
in absolute value by
$6^3\cdot5^6=3375000$, strictly less than that prime.

Permuting the two core orbits and the five other states gives
$2\cdot5!=240$ possible normal-form conjugacies.
The action of the three column-zero band elements and the four
outside actors determines the entire act; column-one band images
are obtained by adding one to the column-zero images.
The complete state-orbit calculation leaves $9139$ representatives
on $40$ of the $43$ targets.  All raw acts and representatives
are retained, so this reduction has an explicit coverage check.
The independent checker
\path{code/d_nakayama21_orbit_coverage.cpp} instead uses all
eleven actor maps as its key and expands all $240$ conjugates
of every retained representative.  Its $2193360$ lookups cover
all $131144$ raw acts, with disjoint retained orbits.
The separate \path{code/d_nakayama21_linkage.py} verifies all
normal-form equations and the literal source-to-category tables;
both saved result files report \texttt{PASS}.

For each representative the converter
\path{code/a_nakayama21_categories.py} builds the literal
$21$-arrow composition table, with identities first, and checks
the monoid and action equations again.  The normalized nerve has
\[
 (\dim C_1,\dim C_2,\dim C_3,\dim C_4)
          =(19,190,1900,19000)
\]
over $\mathbb F_2$.
The exact checker \path{code/a_nakayama21_cups.cpp} verifies
all typed associativity equations, computes a nonboundary
degree-two cocycle $u$, and solves
\[
 \delta w=u\smile u
\]
using the Alexander--Whitney product.  Every one of the $9139$
cases has $H^1=0$, $\dim H^2=1$, $H^3=0$, and such a primitive.
It then checks the primitive on every four-simplex.  In total
this checks $23112531$ typed associativity equations,
$17364100$ cocycle equations, and $173641000$ primitive equations.

The full tables and bit-vector cochains are in
\path{code/outputs/a_nakayama21_cups.jsonl}.
The separate checker \path{code/a_verify_nakayama21_cups.cpp}
imports none of the search or linear-system code: it reconstructs
the normalized chains, computes ranks using the opposite pivot
order, verifies that $u$ is not a coboundary, and evaluates every
stored cocycle and primitive equation directly.
Its full result in
\path{code/outputs/a_nakayama21_cups_independently_verified.json}
is \texttt{PASS} on all $9139$ categories and all the displayed
equations.
Thus the nonzero mod-two square required for $\mathbb CP^2$
is absent in every possible candidate.  This proves the lemma.
\end{proof}

\begin{theorem}[Five outside elements through nine states]
\label{C:eleven-three-row-nine-state}
Let $M$ have a $3\times2$ rectangular-band minimal ideal and
exactly five outside elements.  For every nonempty finite right
$M$-set $S$ with $|S|\leq9$,
\[
 B\mathcal C(M,S)\not\simeq\mathbb CP^2.
\]
\end{theorem}
\begin{proof}
A state fixed by the minimal band gives zero degree-two squares.
Otherwise let $o$ be the number of two-point core orbits and
apply Lemmas~\ref{C:two-column-projective-orbit-kernel} and
\ref{C:two-column-finite-quotient-model}.  The compressed monoid
has the same contracted algebra $A$ and row module $V$ as $M$,
and its column augmentation module $W$ satisfies
\[
 \dim V=2,\qquad \dim W=|S|-2o+1,
 \qquad(\dim\operatorname{Tor}_0,\dim\operatorname{Tor}_1,
        \dim\operatorname{Tor}_2)=(o,0,1),
\]
with zero higher Tor.  The local case is excluded for all $S$
by Corollary~\ref{C:local-five-two-column-act-cones}.
Proposition~\ref{B:five-nonlocal-exceptions} leaves, in the
nonlocal case, only the zero-composite line algebra and the
Nakayama algebra.

For the line algebra the degree-one and degree-two conditions
force $V=S_1\oplus S_3$, by the indecomposable calculation in
Lemma~\ref{B:five-exceptional-targets}, independently of
$\dim W$.  This is a split row module, so
Proposition~\ref{C:split-row-act-cups} applies.
For the Nakayama algebra the same table forces
\[
 V=L,\qquad W=T_2\oplus(e_1A)^{\oplus p}
                    \oplus(e_2A)^{\oplus q},\qquad
 o=p+q.
\]
Comparing dimensions gives
\[
 |S|=4o+p.
\]
If $o=1$, there is one core orbit, hence every band actor is
constant on all of $S$.  The formula permits only $|S|=4$ or
$5$; compression would give a projective-plane monoid of order
$5+3|S|\leq20$, contradicting its exact minimum $23$.

If $o\geq2$ and $|S|\leq9$, the only possibilities are
$(o,p,q,|S|)=(2,0,2,8)$ and $(2,1,1,9)$.
In the first case the compressed monoid has minimal band
$3\times6$ and the Nakayama module pair
$W=T_2\oplus(e_2A)^2$, $V=L$.
The actual-basis quotient argument in the proof of
Proposition~\ref{B:three-by-six-five-outside} makes every
rational degree-two product zero for this pair: its faithful
action quotient has hereditary contracted algebra, and the
degree-two map is an isomorphism.  This part of that proof
does not require the nondegenerate intersection form of its
named product target.  All degree-two classes of the original
act cone pull back from the compressed monoid, so they too
have zero products.  The nine-state case is precisely
Lemma~\ref{C:Nakayama-eleven-nine-state}.
\end{proof}

\begin{corollary}[Small endomorphism monoids through twenty-one arrows]
\label{C:eleven-act-cones21}
An act-cone category $\mathcal C(M,S)$ with $|M|\leq11$ and
$|M|+|S|+1\leq21$ cannot have classifying space
$\mathbb CP^2$.
\end{corollary}
\begin{proof}
Theorem~\ref{C:nine-act-cones} deals with $|M|\leq9$.
A minimal ideal with at most two rows is excluded by
Theorem~\ref{C:two-row-act-cups}.  A minimal ideal with one
column makes the monoid group type; its rational degree-two
cohomology is zero, so relative products in the act cone vanish.
Thus at least three rows and two columns are necessary.
At orders ten and eleven a nontrivial minimal maximal group
would require at least $3\cdot2\cdot2$ ideal elements, already
too many.  The minimal ideal is therefore a band.

The possible shapes are $3\times2$, $4\times2$, $5\times2$,
and $3\times3$, with the identity counted outside the band.
All two-column cases with at most four outside elements are
excluded by Theorem~\ref{C:two-column-four-outside-act-cones}.
The remaining two-column case is $3\times2$ with five outside
elements; here $|S|\leq9$, and the preceding theorem applies.

For the $3\times3$ band the contracted rational algebra has
dimension at most two.  It is either semisimple or is the dual
number algebra $\mathbb Q[\epsilon]/(\epsilon^2)$, as follows
directly from the multiplication of its at most one nonidentity
nonzero basis element.  The two-dimensional row augmentation
module is then either projective over the contracted algebra or
a sum of one-dimensional modules.  The latter case is covered
by Proposition~\ref{C:split-row-act-cups}.  In the former case,
the projective minimal-band ideal gives projective dimension at
most one for that row module over $\mathbb QM$.  The exact
augmentation sequence with projective middle term gives
$\operatorname{pd}_{\mathbb QM}\mathbb Q\leq2$.
The act-cone Tor formula therefore makes its fourth homology
zero.  This exhausts the shapes at these orders.
\end{proof}
